\documentclass{article}
\begin{document}
\section*{Fixed binary augmented native isomorphism}
The source and target are two actual native augmented finite families
for the additive direct image constructed in StandardNativePush.
They are presented by the two tagged four-level flat words,
in occurrence-first and joint order. Address words are assumed to
have distinct raw position tags, not distinct coefficient values.
The augmentation of each side is defined independently. The joint
one first applies the actual joint continuation augmentations and
removes their units inside the relevant branch, then uses the two
actual whole parent augmentations. The multiplicative middle-four is
the coefficient isomorphism, not part of the joint augmentation.
No choice or initiality assumption is imposed on zero-fibre maps.

The forward arrow follows the raw tag between its two Fin positions.
Its whole singleton preimage feeds the middle-four coefficient map.
The inverse and both native inverse equations are proved using the
actual direct-image splits. Preservation is an equality of actual
native morphisms into the terminal shape, not an equality inferred
from the carrier alone. The singleton-preimage chart lemma handles
all outer fibres including empty ones. The grid-to-address lemma
identifies the actual computed chart circuit with the native reorder.
Composing through the chosen raw joint chart proves the demand on
the actual native terminal arrow.

The multiplicative distributivity assumptions are those already
used for the binary chart calculation, plus invertibility of the
multiplicative middle-four. The theorem is fixed at one ordered
presentation and two input occurrences. It is not an isomorphism of
functors on varying native inputs; functoriality of the independent
and dependent operations in the native category, naturality under
arbitrary aggregated arrows, and full all-arity coherence remain.
\end{document}
